\documentclass[11pt]{article}
\usepackage[a4paper,margin=1in]{geometry}
\usepackage{amsmath,amsthm,amssymb}
\usepackage{mathtools}

\newtheorem{theorem}{Theorem}[section]
\newtheorem{proposition}[theorem]{Proposition}
\newtheorem{lemma}[theorem]{Lemma}
\newtheorem{corollary}[theorem]{Corollary}
\theoremstyle{definition}
\newtheorem{definition}[theorem]{Definition}
\newtheorem{remark}[theorem]{Remark}

\newcommand{\R}{\mathbb{R}}
\newcommand{\abs}[1]{\left|#1\right|}
\newcommand{\norm}[1]{\left\|#1\right\|}
\newcommand{\Fr}{\mathcal{A}}
\newcommand{\eps}{\varepsilon}

\title{Saint--Venant Assembly:\\
Sharp $\Fr^{2}$ Stability for the Torsion Energy\\
(Route A, Plan 1, Building Block \textsc{SaintVenantAssembly})}
\author{}
\date{}

\begin{document}
\maketitle

\begin{abstract}
We assemble the sharp Saint--Venant stability inequality
\[
   E(\Omega)-E(B_1)\ge c_{\rm SV}(N,R)\,\Fr(\Omega)^{2},
   \qquad |\Omega|=|B_{1}|,\ \Omega\subset B_{R},
\]
from four upstream building blocks: the quantitative Selection
Principle~(S), the Graph Entry theorem~(G), the Schauder Interpolation
lemma~(I), the Nearly Spherical Closure~(N), glued via the two
quantitative properties of the BDV smooth asymmetry $\alpha$
(\cite{BDV} Lemma~4.2(i)--(ii)) and the suboptimal $\Fr^{4}$ bound of
\cite{BDV} Lemma~5.2. No compactness or contradiction step is used;
every constant is explicit in the upstream constants. The exponent~$2$
is saturated by ellipsoids.
\end{abstract}

\section{Notation, hypotheses, and the BDV black boxes}\label{sec:notation}

\subsection{Sets, energies, and asymmetries}

Throughout $N\ge 2$ and $R\ge 2$. All sets are bounded open subsets of
$\R^{N}$. Write $B_{r}=B_{r}(0)$, $\omega_{N}=|B_{1}|$.

\smallskip\noindent
\emph{Torsion energy.}
$E(\Omega)=\min_{v\in W^{1,2}_{0}(\Omega)}
\bigl(\tfrac12\int|\nabla v|^{2}-\int v\bigr)=-\tfrac12 T(\Omega)$,
where $T(\Omega)=\int_{\Omega}u_{\Omega}$ and $-\Delta u_{\Omega}=1$,
$u_{\Omega}=0$ on $\partial\Omega$.

\smallskip\noindent
\emph{Fraenkel asymmetry.}
$\Fr(\Omega)=\inf\bigl\{|\Omega\Delta B|/|B|:B\text{ ball},|B|=|\Omega|\bigr\}\in[0,2)$.

\smallskip\noindent
\emph{BDV smooth asymmetry.} Following \cite[Def.~4.1]{BDV},
$\alpha(\Omega)=\int_{\Omega\Delta B_{1}(x_{\Omega})}\bigl|1-|x-x_{\Omega}|\bigr|\,dx$,
where $x_{\Omega}$ is the BDV barycenter (\cite{BDV}\,\S3).

\smallskip\noindent
\emph{Torsion quotient.}
$Q_{\alpha}(\Omega):=(E(\Omega)-E(B_{1}))/\alpha(\Omega)$ for
$|\Omega|=|B_{1}|$, $\alpha(\Omega)>0$.

\subsection{The BDV black boxes}

\begin{lemma}[\cite{BDV} Lemma~4.2(i)]\label{lem:BDV42i}
There is $C_{1}=C_{1}(N)>0$ such that for every bounded open
$\Omega\subset\R^{N}$,
\[
   C_{1}\,\alpha(\Omega)\ge\abs{\Omega\Delta B_{1}(x_{\Omega})}^{2}.
\]
In particular, if $|\Omega|=\omega_{N}$,
\begin{equation}\label{eq:42i-Fr}
   \Fr(\Omega)^{2}\le\frac{C_{1}}{\omega_{N}^{2}}\,\alpha(\Omega).
\end{equation}
\end{lemma}

\begin{proof}[Reference]
\cite[Lemma~4.2(i)]{BDV}, proved by one-dimensional annular rearrangement.
Constructive.
\end{proof}

\begin{lemma}[\cite{BDV} Lemma~4.2(ii)]\label{lem:BDV42ii}
For every $R\ge 2$ there is $C_{2}=C_{2}(N,R)>0$ such that for all
bounded open $\Omega_{1},\Omega_{2}\subset B_{R}$,
\[
   \abs{\alpha(\Omega_{1})-\alpha(\Omega_{2})}
   \le C_{2}\,\abs{\Omega_{1}\Delta\Omega_{2}}.
\]
Choosing $\Omega_{2}$ to be the Fraenkel-optimal ball for $\Omega_{1}$
(so $\alpha(\Omega_{2})=0$), for $|\Omega|=\omega_{N}$,
\begin{equation}\label{eq:42ii-Fr}
   \alpha(\Omega)\le C_{2}(N,R)\,\omega_{N}\,\Fr(\Omega).
\end{equation}
\end{lemma}

\begin{remark}[Two-sided $\alpha\leftrightarrow\Fr$ equivalence]
\label{rem:two-sided}
On $\{\Omega\subset B_{R},|\Omega|=\omega_{N}\}$,
\[
   \frac{1}{C_{2}\omega_{N}}\alpha(\Omega)\le\Fr(\Omega)
   \le\frac{\sqrt{C_{1}}}{\omega_{N}}\sqrt{\alpha(\Omega)}.
\]
The right inequality is quadratic (small $\alpha$ implies very small
$\Fr$); the left is linear. Both are constructive.
\end{remark}

\begin{lemma}[\cite{BDV} Lemma~5.2: suboptimal $\Fr^{4}$]\label{lem:BDV52}
There is $C_{8}=C_{8}(N)>0$ such that for every bounded open
$\Omega\subset\R^{N}$ with $0<|\Omega|<\infty$,
\[
   E(\Omega)|\Omega|^{-(N+2)/N}-E(B_{1})|B_{1}|^{-(N+2)/N}
   \ge\frac{1}{C_{8}}\Fr(\Omega)^{4}.
\]
For $|\Omega|=\omega_{N}$,
\begin{equation}\label{eq:52-fixedvol}
   E(\Omega)-E(B_{1})\ge\omega_{N}^{(N+2)/N}\,C_{8}^{-1}\,\Fr(\Omega)^{4}.
\end{equation}
\end{lemma}

This is the Fusco--Maggi--Pratelli-type rearrangement bound (see
\cite{BDV}~p.28, \cite{FMP}); rearrangement + P\'olya--Szeg\H{o}, no
contradiction. Absorbing $\omega_{N}^{(N+2)/N}$ into the constant
($C_{8}':=C_{8}/\omega_{N}^{(N+2)/N}$),
\begin{equation}\label{eq:52-clean}
   E(\Omega)-E(B_{1})\ge(C_{8}')^{-1}\,\Fr(\Omega)^{4}
   \qquad(|\Omega|=\omega_{N}).
\end{equation}

\subsection{The four upstream blocks}

\begin{theorem}[(S) Quantitative Selection]\label{thm:SEL}
There exist $C_{\rm sel}(N,R)>0$, $q_{\rm sel}(N,R)>0$,
$\eps_{\rm sel}(N,R)>0$ such that for every open
$\Omega_{0}\subset B_{R}$ with $|\Omega_{0}|=\omega_{N}$,
$\alpha(\Omega_{0})\le\eps_{\rm sel}$, $Q_{\alpha}(\Omega_{0})<q_{\rm sel}$,
there exists a volume-normalized $\widetilde\Omega\subset B_{R}$ with
$|\widetilde\Omega|=\omega_{N}$ satisfying:
(S1) $\widetilde\Omega$ has $C^{1,\gamma}$ boundary, inheriting BDV
regularity;
(S2) $\tfrac12\alpha(\Omega_{0})\le\alpha(\widetilde\Omega)\le\tfrac32\alpha(\Omega_{0})$;
(S3) $Q_{\alpha}(\widetilde\Omega)\le 2C_{\rm sel}Q_{\alpha}(\Omega_{0})$.
\end{theorem}

\begin{theorem}[(G) Graph Entry]\label{thm:GE}
There is $\alpha_{\rm graph}(N,R)>0$, $C_{H}''(N,R)>0$ such that every
$\widetilde\Omega$ as in (S) with
$\alpha(\widetilde\Omega)\le\alpha_{\rm graph}$ is a $C^{1,\gamma}$ normal
graph over $\partial B_{1}$:
$\partial\widetilde\Omega=\{(1+g(\theta))\theta\}$ with
$\norm{g}_{L^{\infty}}\le C_{H}''\alpha(\widetilde\Omega)^{1/(2N)}$.
\end{theorem}

\begin{lemma}[(I) Schauder Interpolation]\label{lem:SI}
There is $\alpha_{\rm sph}(N,R)>0$, $\delta_{\rm sph}(N,\gamma_{0})>0$
such that whenever (G) applies and $\alpha(\widetilde\Omega)\le\alpha_{\rm sph}$,
$\norm{g}_{C^{2,\gamma_{0}}}\le\delta_{\rm sph}$.
\end{lemma}

\begin{theorem}[(N) Nearly Spherical Closure]\label{thm:NS}
There is $c_{*}=c_{*}(N)>0$ such that every nearly spherical
$\widetilde\Omega=\{(1+g(\theta))\theta\}$ with
$\norm{g}_{C^{2,\gamma_{0}}}\le\delta_{\rm sph}$,
$|\widetilde\Omega|=\omega_{N}$, $x_{\widetilde\Omega}=0$, satisfies
$Q_{\alpha}(\widetilde\Omega)\ge c_{*}(N)$.
\end{theorem}

\section{Small-asymmetry regime: linear-in-$\alpha$ bound}\label{sec:small}

\begin{equation}\label{eq:alpha-small-def}
   \alpha_{\rm small}(N,R)
   :=\tfrac12\min\{\eps_{\rm sel},\alpha_{\rm graph},\alpha_{\rm sph}\}>0.
\end{equation}

\begin{theorem}[Small-$\alpha$ linear bound]\label{thm:small-alpha}
With $q_{*}(N,R):=c_{*}(N)/(2C_{\rm sel}(N,R))$, every open
$\Omega_{0}\subset B_{R}$ with $|\Omega_{0}|=\omega_{N}$ and
$\alpha(\Omega_{0})\le\alpha_{\rm small}(N,R)$ satisfies
\begin{equation}\label{eq:linear-alpha}
   E(\Omega_{0})-E(B_{1})\ge q_{*}(N,R)\,\alpha(\Omega_{0}).
\end{equation}
\end{theorem}

\begin{proof}
If $Q_{\alpha}(\Omega_{0})\ge q_{\rm sel}$ then trivially
$E-E(B_{1})\ge q_{\rm sel}\alpha\ge q_{*}\alpha$ (after shrinking
$q_{\rm sel}$ if needed so $q_{*}\le q_{\rm sel}$).

Otherwise $Q_{\alpha}(\Omega_{0})<q_{\rm sel}$. Apply (S):
$|\widetilde\Omega|=\omega_{N}$, $\alpha(\widetilde\Omega)\le\tfrac32\alpha_{\rm small}<\alpha_{\rm graph}$.
Apply (G): graph representation, $\norm{g}_{L^{\infty}}\le C_{H}''\alpha^{1/(2N)}$.
Apply (I): $\norm{g}_{C^{2,\gamma_{0}}}\le\delta_{\rm sph}$ (since
$\alpha(\widetilde\Omega)\le\tfrac32\alpha_{\rm small}<\alpha_{\rm sph}$).
Apply (N): $Q_{\alpha}(\widetilde\Omega)\ge c_{*}(N)$.
By (S3), $c_{*}\le Q_{\alpha}(\widetilde\Omega)\le 2C_{\rm sel}Q_{\alpha}(\Omega_{0})$,
i.e., $Q_{\alpha}(\Omega_{0})\ge q_{*}$.
\end{proof}

\begin{remark}
No limit, no compactness, no contradiction. Every step is a quantitative
implication.
\end{remark}

\section{Conversion $\alpha\Rightarrow\Fr^{2}$}\label{sec:convert}

\begin{corollary}[Sharp small-$\Fr$ quadratic bound]\label{cor:smallFr}
For every open $\Omega\subset B_{R}$ with $|\Omega|=\omega_{N}$ and
$\alpha(\Omega)\le\alpha_{\rm small}(N,R)$,
\begin{equation}\label{eq:small-fr}
   E(\Omega)-E(B_{1})\ge\kappa_{\rm near}(N,R)\,\Fr(\Omega)^{2},
   \qquad
   \kappa_{\rm near}(N,R):=\frac{q_{*}(N,R)\,\omega_{N}^{2}}{C_{1}(N)}.
\end{equation}
\end{corollary}

\begin{proof}
\eqref{eq:linear-alpha} gives $E-E(B_{1})\ge q_{*}\alpha\ge
q_{*}(\omega_{N}^{2}/C_{1})\Fr^{2}$ by~\eqref{eq:42i-Fr}.
\end{proof}

\begin{lemma}[$\Fr$-version of the small regime]\label{lem:Fr-small}
Set
\begin{equation}\label{eq:a-near-def}
   a_{\rm near}(N,R):=\frac{\alpha_{\rm small}(N,R)}{C_{2}(N,R)\,\omega_{N}}>0.
\end{equation}
Every $\Omega\subset B_{R}$ with $|\Omega|=\omega_{N}$ and
$\Fr(\Omega)\le a_{\rm near}(N,R)$ satisfies
$\alpha(\Omega)\le\alpha_{\rm small}$ and hence
$E(\Omega)-E(B_{1})\ge\kappa_{\rm near}\,\Fr(\Omega)^{2}$.
\end{lemma}

\begin{proof}
\eqref{eq:42ii-Fr} gives $\alpha(\Omega)\le C_{2}\omega_{N}\Fr(\Omega)\le
C_{2}\omega_{N}a_{\rm near}=\alpha_{\rm small}$.
\end{proof}

\begin{remark}[Why the two-sided equivalence matters]
The conversion uses \emph{both} directions of BDV Lemma~4.2:
(i) gives the sharp exponent~2 in~\eqref{eq:small-fr};
(ii), applied as above, gives a positive $\Fr$-threshold $a_{\rm near}$
inside which the small-$\alpha$ analysis applies. Without~(ii) the
small-$\alpha$ chain could not be invoked at arbitrary small $\Fr$, and
the gluing would only give the suboptimal $\Fr^{4}$ from BDV Lemma~5.2.
\end{remark}

\section{Far-asymmetry regime}\label{sec:far}

\begin{lemma}[Far-$\Fr$ quadratic bound]\label{lem:far}
For $a>0$ and every open $\Omega\subset B_{R}$ with $|\Omega|=\omega_{N}$
and $\Fr(\Omega)\ge a$,
\begin{equation}\label{eq:far}
   E(\Omega)-E(B_{1})\ge(C_{8}')^{-1}\,\Fr(\Omega)^{4}
   \ge\frac{a^{2}}{C_{8}'(N)}\,\Fr(\Omega)^{2}.
\end{equation}
\end{lemma}

\begin{proof}
\eqref{eq:52-clean} gives the first inequality; $\Fr^{4}=\Fr^{2}\Fr^{2}\ge
a^{2}\Fr^{2}$ gives the second.
\end{proof}

\section{Gluing}\label{sec:glue}

\begin{theorem}[Sharp Saint--Venant, single-statement form]\label{thm:SV}
Define
\begin{align}
   a_{\rm far}(N,R)&:=a_{\rm near}(N,R)=\frac{\alpha_{\rm small}(N,R)}{C_{2}(N,R)\,\omega_{N}},
   \label{eq:afar-def}\\
   c_{\rm SV}(N,R)&:=\min\!\Bigl\{\kappa_{\rm near}(N,R),\frac{a_{\rm far}(N,R)^{2}}{C_{8}'(N)}\Bigr\}>0.
   \label{eq:cSV-def}
\end{align}
For every open $\Omega\subset B_{R}$ with $|\Omega|=\omega_{N}$,
\begin{equation}\label{eq:SV-final}
   E(\Omega)-E(B_{1})\ge c_{\rm SV}(N,R)\,\Fr(\Omega)^{2}.
\end{equation}
\end{theorem}

\begin{proof}
\emph{Case A} ($\Fr\le a_{\rm near}$): Lemma~\ref{lem:Fr-small} gives
$E-E(B_{1})\ge\kappa_{\rm near}\Fr^{2}\ge c_{\rm SV}\Fr^{2}$.
\emph{Case B} ($\Fr\ge a_{\rm far}$): Lemma~\ref{lem:far} with
$a=a_{\rm far}$ gives $E-E(B_{1})\ge(a_{\rm far}^{2}/C_{8}')\Fr^{2}\ge c_{\rm SV}\Fr^{2}$.
Since $a_{\rm near}=a_{\rm far}$ the two cases cover $[0,2)$.
\end{proof}

\begin{remark}[Sharpness]
Equality in~\eqref{eq:SV-final} is asymptotically attained along small
volume-preserving ellipsoidal perturbations of $B_{1}$: in that family
$\Fr\sim\norm{g}_{L^{1}}$, $E-E(B_{1})\sim\norm{g}_{L^{2}}^{2}$ to
second order (\cite{BDV} Introduction, p.2). The exponent~$2$ is sharp.
\end{remark}

\section{Final statement}\label{sec:thm}

\begin{theorem}[\textsc{SaintVenantAssembly}]\label{thm:SVA-main}
Let $N\ge 2$, $R\ge 2$. There exists $c_{\rm SV}(N,R)>0$, explicit in
$C_{\rm sel}(N,R), c_{*}(N), \alpha_{\rm small}(N,R), C_{1}(N), C_{2}(N,R),
C_{8}'(N)$ via~\eqref{eq:small-fr}, \eqref{eq:a-near-def},
\eqref{eq:cSV-def}, such that for every open $\Omega\subset B_{R}$ with
$|\Omega|=|B_{1}|$,
\[
   \boxed{\;E(\Omega)-E(B_{1})\ge c_{\rm SV}(N,R)\,\Fr(\Omega)^{2}.\;}
\]
The exponent~2 is sharp (saturated by ellipsoidal perturbations).
$c_{\rm SV}(N,R)$ is extracted from upstream constants without any
compactness, contradiction or limit argument. \qed
\end{theorem}

\section{Summary of constants}\label{sec:summary}

\renewcommand{\arraystretch}{1.25}
\begin{center}
\begin{tabular}{|l|l|l|}\hline
Symbol & Source & Meaning \\\hline
$\omega_{N}$ & $|B_{1}|$ & ball volume \\
$C_{1}(N)$ & BDV L4.2(i) & $C_{1}\alpha\ge|\Omega\Delta B|^{2}$ \\
$C_{2}(N,R)$ & BDV L4.2(ii) & Lipschitz $\alpha$ in $L^{1}$ \\
$C_{8}'(N)$ & BDV L5.2 (rescaled) & $E-E(B_{1})\ge(C_{8}')^{-1}\Fr^{4}$ \\
$C_{\rm sel}(N,R),\eps_{\rm sel},q_{\rm sel}$ & block (S) & selection \\
$\alpha_{\rm graph}(N,R),C_{H}''$ & block (G) & graph entry \\
$\alpha_{\rm sph}(N,R),\delta_{\rm sph}$ & blocks (I), (N) & Schauder/nearly sph.\ \\
$c_{*}(N)$ & block (N) & nearly-spherical floor \\\hline
$\alpha_{\rm small}$ & \eqref{eq:alpha-small-def} & $\tfrac12\min\{\eps_{\rm sel},\alpha_{\rm graph},\alpha_{\rm sph}\}$ \\
$q_{*}=c_{*}/(2C_{\rm sel})$ & Thm~\ref{thm:small-alpha} & small-$\alpha$ slope \\
$\kappa_{\rm near}=q_{*}\omega_{N}^{2}/C_{1}$ & \eqref{eq:small-fr} & small-$\Fr$ const. \\
$a_{\rm near}=a_{\rm far}=\alpha_{\rm small}/(C_{2}\omega_{N})$ & \eqref{eq:a-near-def} & glue threshold \\
$c_{\rm SV}(N,R)$ & \eqref{eq:cSV-def} & $\min\{\kappa_{\rm near},a_{\rm far}^{2}/C_{8}'\}$ \\\hline
\end{tabular}
\end{center}

\subsection*{Where compactness would have appeared}

\begin{itemize}
\item BDV Theorem~4.3 (linear-in-$\alpha$ inequality) is proved
in~\cite{BDV} by contradiction. We replaced it with
Theorem~\ref{thm:small-alpha}, deriving $q_{*}=c_{*}/(2C_{\rm sel})$
constructively from (S)+(G)+(I)+(N).
\item BDV Lemma~5.2 ($\Fr^{4}$) is used as a black box; its proof is
rearrangement-based.
\item The conversion uses BDV Lemma~4.2(i)+(ii), both constructive.
\item The gluing is arithmetic.
\end{itemize}

No compactness, no minimizing sequence, no limit set, no contradiction
hypothesis appears at any stage.

\begin{thebibliography}{9}
\bibitem{BDV} L.~Brasco, G.~De~Philippis, B.~Velichkov,
\emph{Faber--Krahn inequalities in sharp quantitative form},
Duke Math.\ J.\ \textbf{164} (2015), 1777--1831 (arXiv:1306.0392).
\bibitem{FMP} N.~Fusco, F.~Maggi, A.~Pratelli,
\emph{The sharp quantitative Sobolev inequality for functions of bounded
variation}, J.~Funct.~Anal.\ \textbf{244} (2007), 315--341.
\end{thebibliography}

\end{document}
